\makeatletter
\def\input@path{{./}{../}{../../}{../../../}{../../../../}{preamble/}{../preamble/}{../../preamble/}{../../../preamble/}}
\makeatother

\input{preamble_exams.tex}
\usepackage{needspace}
\usepackage{titling}
\setlength{\droptitle}{-2em}
\pretitle{\par\begingroup\centering\LARGE}
\posttitle{\par\endgroup}
\preauthor{}\postauthor{}
\predate{}\postdate{}
\maketitlehookd{\vspace{-1em}}
\linespread{1.12}
\setlength{\parskip}{5pt}
\newgeometry{margin=0.65in,top=0.55in,bottom=0.65in}

\title{Grade 11 Term 2 Home Practice -- C8\\\large Student's \(t\) interval for a mean\\[0.25em]
	\normalsize Global Economics -- 2026--2027}
\author{}
\date{}

\begin{document}
	\maketitle

	\setcounter{tocdepth}{1}
	\setcounter{secnumdepth}{2}
	\phantomsection
	\label{toc}
	\tableofcontents

	\section*{Evaluated criteria}
	\begin{enumerate}[label=\textbf{C\arabic*},start=8]
		\item Constructs confidence intervals for a population mean or population proportion with unknown variance using the Student's t-distribution, interpreting the result in finance and economics contexts.
	\end{enumerate}

	\textbf{How to practise.} Try each task before opening its blue solution link. The arrows return to the contents, problem list or current problem.

	\section{Worked Solution Examples}
	\label{sec:worked-solution-examples}

	\SectionProblemsTableOfContents{
		\SectionProblemLink{prob:c8-1}{Loading time and dispatch costs (C8)} \\
		\SectionProblemLink{prob:c8-2}{Customs processing for an export route (C8)}
	}

	% ----------------------------------------------------------------
	\subsection{Loading time and dispatch costs (C8)}
	\label{prob:c8-1}

	\noindent
	\begin{minipage}[t]{0.53\textwidth}
	\vspace{0pt}
	A logistics centre uses loading time to plan labour costs in a new dispatch lane. From a roster of 400 eligible loads for the month, it draws a simple random sample of 12 distinct loads. The times have a roughly symmetric dotplot with no unusual values. The sample gives $\bar x=84.60$ minutes and $s=6.80$ minutes. The population standard deviation is unknown.
	\end{minipage}
	\hfill
	\begin{minipage}[t]{0.43\textwidth}
	\vspace{0pt}
	\centering\small
	Student's $t$ cutoffs by degrees of freedom and central area.\par\smallskip
	\begin{tabular}{@{}rcc@{}}
	\toprule
	$df$ & 90\% & 95\% \\
	\midrule
	11 & 1.796 & 2.201 \\
	12 & 1.782 & 2.179 \\
	\bottomrule
	\end{tabular}
	\end{minipage}

	\QuestionsTasks[sec:worked-solution-examples]{
		{8}/{sol:c8-1-a}/
		{Define $\mu$ and justify a one-sample Student's $t$ interval. Use the table to find the cutoff.},
		{8}/{sol:c8-1-b}/
		{Compute the standard error, margin of error and endpoints. Explain the result for the centre's labour planning.}
	}

	\clearpage

	\Needspace{5\baselineskip}
	\phantomsection
	\label{sol:c8-1-a}
	\subsubsection*{C8 -- Justify the method and read the t cutoff}

	The parameter $\mu$ is the mean loading time of the 400 eligible loads. Sampling was random and $12/400=3\%$, supporting approximate independence. With a small sample, the roughly symmetric dotplot and lack of unusual values support the shape assumption. We estimate the unknown population SD using $s$, so the standardized statistic uses a Student's $t$ reference distribution under the model. As in the Term 1 table work, read the two tails together: $df=n-1=11$, and a central 95\% leaves 0.025 in each tail, giving $t^*=2.201$.

	\secInicio[prob:c8-1][sec:worked-solution-examples]

	\Needspace{5\baselineskip}
	\phantomsection
	\label{sol:c8-1-b}
	\subsubsection*{C8 -- Construct and interpret the t interval}

	\[
	\widehat{SE}=\frac{s}{\sqrt n}=\frac{6.80}{\sqrt{12}}\approx1.9630\text{ min},
	\qquad E=2.201\left(\frac{6.80}{\sqrt{12}}\right)\approx4.3205\text{ min}.
	\]
	The 95\% interval is $84.60\pm E$, or approximately $[80.28,88.92]$ minutes. We are 95\% confident the mean loading time for the eligible loads is in this range. This gives a range for the average loading time used in labour budgeting. Individual loads can take less or more time than these endpoints.

	\secInicio[prob:c8-1][sec:worked-solution-examples]

	\clearpage

	% ----------------------------------------------------------------
	\subsection{Customs processing for an export route (C8)}
	\label{prob:c8-2}

	An exporter draws a simple random sample without replacement of 18 shipments from 600 using one route last season. Customs processing times have a roughly symmetric sample distribution with no strong outliers. The mean is 42.30 hours and the sample SD is 5.10 hours. The population SD is unknown. For 90\% central area, use $t^*=1.740$ with 17 degrees of freedom and $z^*=1.645$ for the standard normal reference.

	\QuestionsTasks[sec:worked-solution-examples]{
		{8}/{sol:c8-2-a}/
		{Identify the population mean, justify which critical value applies, and calculate a two-sided 90\% interval.},
		{8}/{sol:c8-2-b}/
		{State the result in context and assess the claim that the route's population mean processing time is below 40 hours.}
	}

	\Needspace{5\baselineskip}
	\phantomsection
	\label{sol:c8-2-a}
	\subsubsection*{C8 -- Choose the cutoff and calculate the interval}

	The target is $\mu$, the mean customs processing time of the 600 shipments on this route last season. Random selection, $18/600=3\%$, and the stated shape support the method. The procedure estimates the unknown $\sigma$ using $s$, so use the $t$ cutoff with $df=18-1=17$.
	\[
	\widehat{SE}=\frac{5.10}{\sqrt{18}}\approx1.2021\text{ h},\qquad
	E=1.740\left(\frac{5.10}{\sqrt{18}}\right)\approx2.0916\text{ h}.
	\]
	The interval is $42.30\pm E$, or approximately $[40.21,44.39]$ hours.

	\secInicio[prob:c8-2][sec:worked-solution-examples]

	\Needspace{5\baselineskip}
	\phantomsection
	\label{sol:c8-2-b}
	\subsubsection*{C8 -- Interpret the interval for the exporter's schedule}

	We are 90\% confident that the mean customs processing time for the 600 shipments last season was about 40.21 to 44.39 hours. All values in the interval are above 40 hours, so the data give evidence that the season's population mean exceeds 40 hours. A schedule for individual shipments also needs to account for variability around that mean.

	\secInicio[prob:c8-2][sec:worked-solution-examples]


\end{document}
